% !TeX program = lualatex
% Example deck for beamerthemefulda.sty. All content is placeholder text.
\documentclass[aspectratio=169, 11pt, english]{beamer}

\usepackage[english]{babel}
\usepackage{booktabs}

% Shows hs-fulda-logo.png from this folder. Other file: logo=<file>. No logo: nologo.
\usetheme{fulda}

\title{Presentation title}
\subtitle{A subtitle or guiding question}
\author{First Author \and Second Author}
\institute{Department or module, lecturer}
\date{Date of the presentation}

\begin{document}

\begin{frame}[plain]
  \titlepage
\end{frame}

\begin{frame}{Agenda}
  \tableofcontents
\end{frame}

\section{Lists and emphasis}

\begin{frame}{Lists and emphasis}
  \framesubtitle{A frame subtitle is optional}

  \begin{itemize}
    \item One statement per bullet
    \item \alert{Emphasis} with \texttt{\textbackslash alert}
    \begin{itemize}
      \item Second level
    \end{itemize}
  \end{itemize}

  \begin{enumerate}
    \item First step
    \item Second step
  \end{enumerate}
\end{frame}

\section{Columns and blocks}

\begin{frame}{Columns and blocks}
  \begin{columns}[T, onlytextwidth]
    \begin{column}{0.48\textwidth}
      Text and blocks can sit side by side.

      The content starts at the same height on every slide.

      \begin{alertblock}{Alert block}
        For warnings and limits.
      \end{alertblock}
    \end{column}
    \begin{column}{0.48\textwidth}
      \begin{block}{Block}
        A block highlights a key statement.
      \end{block}
      \begin{exampleblock}{Example block}
        For cases and exercises.
      \end{exampleblock}
    \end{column}
  \end{columns}
\end{frame}

\section{Quotes, tables and sources}

\begin{frame}{Quotes, tables and sources}
  \pullquote[Author (Year), p.~34]{“A verbatim quote that should stand out on the slide.”}

  \begin{center}
    \begin{tabular}{lrr}
      \toprule
      Condition & Before & After \\
      \midrule
      Group A & 0.0 & 0.0 \\
      Group B & 0.0 & 0.0 \\
      \bottomrule
    \end{tabular}
  \end{center}

  \sourceline{placeholder values, own illustration}
\end{frame}

\section*{References}

\begin{frame}{References}
  \begin{thebibliography}{9}
    \bibitem{author-a}
      Author, A. (Year). \emph{Title of the book}. Publisher.
    \bibitem{author-b}
      Author, B., \& Author, C. (Year). Title of the article. \emph{Journal Name}, 1(2), 34--56.
  \end{thebibliography}
\end{frame}

\end{document}
